% ЗГЕНЕРОВАНО scripts/build_format_v2_chapter.py — не редагувати вручну.
% Джерело: editorial/cards/9.*.json, format-v2/rozdil-9.json, paraleli.json.

\naukchapter{9}{8}{Глава 9. Відречення}
\naukintro{Девʼята глава — вісім віршів про відречення, і жодного разу воно не постає як зусилля. Спершу пари протилежностей утихають самі (9.1), далі втихає бажання жити (9.2), далі річ виявляється без суті (9.3). Найгостріший вірш — 9.8: сансара названа просто усталеними схильностями, і тоді відречення зводиться до одного — відпустити їх.}
\vsec{9.1}
\vdeva{कृताकृते च द्वन्द्वानि कदा शान्तानि कस्य वा । एवं ज्ञात्वेह निर्वेदाद् भव त्यागपरो ऽव्रती ॥}
\viast{kṛtākṛte ca dvandvāni kadā śāntāni kasya vā | evaṃ jñātveha nirvedād bhava tyāgaparo 'vratī ||}
\vtranslit{крітакріте ча двандвані када шантані кася ва | евам джнятвеха нірведад бгава тяґапаро ’враті ||}
\vperek{Зроблене й незроблене, пари протилежностей — коли вони втихли і чиї вони? Так знаючи, будь тут, у відреченні, схильним до відкидання — і без обітниць.}
\vcom{Построкове зіставлення: редакторська чернетка, не рецензовано.}
\vcom{Рядок а}
\vdeva{कृताकृते च द्वन्द्वानि कदा शान्तानि कस्य वा}
\viast{kṛtākṛte ca dvandvāni kadā śāntāni kasya vā}
\vtranslit{крітакріте ча двандвані када шантані кася ва}
\vperek{Зроблене й незроблене, пари протилежностей — коли вони втихли і чиї вони?}
\vcom{Питальні слова kadā («коли?») і kasya vā («чиї?») завершують перший друкований піврядок. dvandvāni — «пари протилежностей»; kṛtākṛte потребує окремої морфологічної звірки, тому українське «зроблене й незроблене» лишається чинним робочим відповідником, а не остаточно встановленим граматичним рішенням.}
\vcom{Рядок б}
\vdeva{एवं ज्ञात्वेह निर्वेदाद् भव त्यागपरो ऽव्रती}
\viast{evaṃ jñātveha nirvedād bhava tyāgaparo 'vratī}
\vtranslit{евам джнятвеха нірведад бгава тяґапаро ’враті}
\vperek{Так знаючи, будь тут, у відреченні, схильним до відкидання — і без обітниць.}
\vcom{evaṃ jñātvā («так пізнавши») передує наказовому bhava («будь»); nirvedād вказує на підставу відречення. tyāgaparaḥ («схильний до відкидання») поєднано з avratī («без обітниць»): внутрішнє відречення не перетворюється на припис нової обітниці. GRETIL і SARIT по-різному позначають сегментацію nirvedād bhava / nirvedādbhava та 'vratī / +avratī; це не доказ незалежного рукописного різночитання.}
\vcom{`avratī` («без обітниць») — не дрібниця: текст водночас радить `tyāga` (відкидання) і знімає `vrata` (обітницю). Тобто відкидання не повинно стати новою дисципліною — інакше воно перетвориться на те, від чого відкидають. Питання `kasya vā` («і чиї ж вони?») залишене без відповіді навмисно.}
\vpar{Українська паралель (наша інтерпретація)}{Олександр Потебня, «Думка й мова» (1862).}{Наша паралель: Потебня теж питає, чиї насправді ті пари, якими ми міряємо світ, — і відповідь у нього не в речах, а в тому, хто міряє. Наше порівняння — типологічне.}
\vfig{ilustratsii/9-1-neti.pdf}{Пари протилежностей втихли: лишається відречення без обітниць}{Схема наша, не історична ілюстрація.}

\vsec{9.2}
\vdeva{कस्यापि तात धन्यस्य लोकचेष्टावलोकनात् । जीवितेच्छा बुभुक्षा च बुभुत्सोपशमं गताः ॥}
\viast{kasyāpi tāta dhanyasya lokaceṣṭāvalokanāt | jīvitecchā bubhukṣā ca bubhutsopaśamaṃ gatāḥ ||}
\vtranslit{касяпі тата дганяся локачештавалоканат | джівітечча бубгукша ча бубгутсопашамам ґатах ||}
\vperek{У якогось щасливого, любий, від споглядання того, як поводиться світ, бажання жити, голод і жага пізнавати дійшли до втишення.}
\vcom{Построкове зіставлення: редакторська чернетка, не рецензовано.}
\vcom{Рядок а}
\vdeva{कस्यापि तात धन्यस्य लोकचेष्टावलोकनात्}
\viast{kasyāpi tāta dhanyasya lokaceṣṭāvalokanāt}
\vtranslit{касяпі тата дганяся локачештавалоканат}
\vperek{У якогось щасливого, любий, від споглядання того, як поводиться світ,}
\vcom{Родовий зворот kasyāpi ... dhanyasya («якогось щасливого») окреслює особу, а tāta — звертання до співрозмовника («любий»). Аблатив lokaceṣṭāvalokanāt пояснює підставу: від споглядання руху й учинків світу. Це повідомлення про досвід, а не наказ споглядати за певним правилом.}
\vcom{Рядок б}
\vdeva{जीवितेच्छा बुभुक्षा च बुभुत्सोपशमं गताः}
\viast{jīvitecchā bubhukṣā ca bubhutsopaśamaṃ gatāḥ}
\vtranslit{джівітечча бубгукша ча бубгутсопашамам ґатах}
\vperek{бажання жити, голод і жага пізнавати дійшли до втишення.}
\vcom{jīvitecchā («бажання жити») та bubhukṣā («голод») названі разом; у третьому понятті йдеться про жагу пізнавати. SARIT має bubhutsopaśamaṃ gatāḥ, а GRETIL/TextGrid, за задокументованою звіркою Issue \#227, — bubhutsopaśamaḥ gatāḥ: різниця -ṃ / -ḥ відкрита й змінює можливий синтаксичний розбір. Не вважати електронні читання рукописною колацією. Втишення бажань у цьому контексті не дорівнює заклику до бездіяльності.}
\vcom{Три подібні іменники стоять поруч: `jīvitecchā` (бажання жити), `bubhukṣā` (голод) і `bubhutsā` (жага пізнавати) — останнє те саме слово, що в 10.5, і словник має бути спільним. Причина названа одна й зовнішня: **споглядання** того, як поводиться світ, а не зусилля над собою. `tāta` — ласкаве звертання, не «батько».}
\vpar{Українська паралель (наша інтерпретація)}{Леся Українка, «Оргія» (1913).}{Наша паралель: у Лесі Українки бажання жити теж утихає не від зусилля, а від побаченого. Наше порівняння — за способом бачити.}
\vfig{ilustratsii/9-2-flame.pdf}{Полумʼя, що втихло: бажання жити й жага пізнавати}{Схема наша, не історична ілюстрація.}

\vsec{9.3}
\vdeva{अनित्यं सर्वम् एवेदं तापत्रितयदूषितम् । असारं निन्दितं हेयम् इति निश्चित्य शाम्यति ॥}
\viast{anityaṃ sarvam evedaṃ tāpatritayadūṣitam | asāraṃ ninditaṃ heyam iti niścitya śāmyati ||}
\vtranslit{анітям сарвам еведам тапатрітаядушітам | асарам ніндітам хеям іті нішчітя шам'яті ||}
\vperek{Усе це нестале, зіпсоване троїстою мукою, без суті, осуджене, гідне відкинення — так вирішивши, людина втихає.}
\vcom{Построкове зіставлення: редакторська чернетка, не рецензовано.}
\vcom{Рядок а}
\vdeva{अनित्यं सर्वम् एवेदं तापत्रितयदूषितम्}
\viast{anityaṃ sarvam evedaṃ tāpatritayadūṣitam}
\vtranslit{анітям сарвам еведам тапатрітаядушітам}
\vperek{Усе це нестале, зіпсоване троїстою мукою,}
\vcom{anityaṃ («нестале») характеризує sarvam eva idam («усе це»); tāpa-tritaya-dūṣitam — «уражене троїстою мукою». Сам цей піврядок не називає окремо трьох її різновидів: конкретний перелік належить до пояснення традиції, а не до буквального санскритського рядка. Українське «усе це» зберігає вказівний наголос idam.}
\vcom{Рядок б}
\vdeva{असारं निन्दितं हेयम् इति निश्चित्य शाम्यति}
\viast{asāraṃ ninditaṃ heyam iti niścitya śāmyati}
\vtranslit{асарам ніндітам хеям іті нішчітя шам'яті}
\vperek{без суті, осуджене, гідне відкинення — так вирішивши, людина втихає.}
\vcom{asāraṃ («без суті»), ninditaṃ («осуджене»), heyam («те, що належить відкинути») продовжують низку оцінок мінливого. iti niścitya — «так визначивши / дійшовши висновку», після чого śāmyati означає «втихає». Ця синтаксична послідовність описує втишення після розпізнання, а не наказ нехтувати щоденними обов'язками чи здоров'ям.}
\vcom{`tāpa-tritaya` («троїста мука») — усталена формула: мука від себе, від інших і від природи; переклад має лишити її як формулу, а не розтлумачувати в тексті вірша. П'ять означень ідуть підряд, і кожне сильніше за попереднє: від `anitya` (нестале) до `heya` (гідне відкинення). Втишення тут — **наслідок рішення**, а не почуття.}
\vpar{Українська паралель (наша інтерпретація)}{Ольга Кобилянська, «Земля» (1902).}{Наша паралель: у Кобилянської речі, за які тримаються, теж виявляються несталими й без суті. Наше порівняння — типологічне, не твердження про джерело.}
\vfig{ilustratsii/9-3-neti.pdf}{Нестале й троїста мука відкинуті — людина втихає}{Схема наша, не історична ілюстрація.}

\vsec{9.4}
\vdeva{को ऽसौ कालो वयः किं वा यत्र द्वन्द्वानि नो नृणाम् । तान्य् उपेक्ष्य यथाप्राप्तवर्ती सिद्धिम् अवाप्नुयात् ॥}
\viast{ko 'sau kālo vayaḥ kiṃ vā yatra dvandvāni no nṛṇām | tāny upekṣya yathāprāptavartī siddhim avāpnuyāt ||}
\vtranslit{ко ’сау кало ваях кім ва ятра двандвані но нрінам | танй упекшя ятхапраптаварті сіддгім авапнуят ||}
\vperek{Який же той час, вік або що — коли в людей немає пар протилежностей? Знехтувавши їх, хто живе як випало, той досяг би здійснення.}
\vcom{Построкове зіставлення: редакторська чернетка, не рецензовано.}
\vcom{Рядок а}
\vdeva{को ऽसौ कालो वयः किं वा यत्र द्वन्द्वानि नो नृणाम्}
\viast{ko 'sau kālo vayaḥ kiṃ vā yatra dvandvāni no nṛṇām}
\vtranslit{ко ’сау кало ваях кім ва ятра двандвані но нрінам}
\vperek{Який же той час, вік або що — коли в людей немає пар протилежностей?}
\vcom{Початкове ko 'sau kālo («який же той час») та vayaḥ kiṃ vā («вік або що») утворюють риторичне питання. У сполученні yatra dvandvāni no nṛṇām слово no передає заперечення наявності «пар протилежностей» (dvandvāni) серед людей (nṛṇām), а не присвійне «наших». Питання не обіцяє, що певний вік звільнить людину від протилежностей.}
\vcom{Рядок б}
\vdeva{तान्य् उपेक्ष्य यथाप्राप्तवर्ती सिद्धिम् अवाप्नुयात्}
\viast{tāny upekṣya yathāprāptavartī siddhim avāpnuyāt}
\vtranslit{танй упекшя ятхапраптаварті сіддгім авапнуят}
\vperek{Знехтувавши їх, хто живе як випало, той досяг би здійснення.}
\vcom{tāny upekṣya («не зважаючи на них / відкинувши їх») відсилає до попередніх dvandvāni. yathāprāptavartī буквально окреслює людину, яка живе відповідно до того, що трапляється; це не твердження, ніби обов'язки зникають. siddhim avāpnuyāt містить потенційно-умовну форму avāpnuyāt («досяг би»): переклад не повинен перетворювати її на безумовну обіцянку.}
\vcom{Питання риторичне й знецінювальне: не існує такого віку, коли пар немає, — отже, чекати на них марно. `yathāprāpta-vartī` («той, хто живе як випало») — той самий принцип, що `yadṛcchā` у 3.14 і 4.4, і в перекладі він має звучати так само. `avāpnuyāt` — умовний спосіб: «досяг би», а не «досягне»; це не обіцянка.}
\vpar{Українська паралель (ПОТРЕБУЄ ПЕРЕВІРКИ)}{українська філософська лірика, не визначено.}{ПОТРЕБУЄ ПЕРЕВІРКИ: думка, що жити як випало — це і є здійснення, близька до української філософської лірики, але конкретного твору ми не звірили.}
\vfig{ilustratsii/9-4-balance.pdf}{Дві шальки: пари протилежностей і життя як випало}{Схема наша, не історична ілюстрація.}

\vsec{9.5}
\vdeva{नाना मतं महर्षीणां साधूनां योगिनां तथा । दृष्ट्वा निर्वेदम् आपन्नः को न शाम्यति मानवः ॥}
\viast{nānā mataṃ maharṣīṇāṃ sādhūnāṃ yogināṃ tathā | dṛṣṭvā nirvedam āpannaḥ ko na śāmyati mānavaḥ ||}
\vtranslit{нана матам махаршінам садгунам йоґінам татха | дріштва нірведам апаннах ко на шам'яті манавах ||}
\vperek{Побачивши, скільки різних поглядів мають великі мудреці, праведні й йоги, — хто, дійшовши відречення, не втихне?}
\vcom{Построкове зіставлення: редакторська чернетка, не рецензовано.}
\vcom{Рядок а}
\vdeva{नाना मतं महर्षीणां साधूनां योगिनां तथा}
\viast{nānā mataṃ maharṣīṇāṃ sādhūnāṃ yogināṃ tathā}
\vtranslit{нана матам махаршінам садгунам йоґінам татха}
\vperek{Побачивши, скільки різних поглядів мають великі мудреці, праведні й йоги, —}
\vcom{nānā mataṃ («різноманітні погляди») стосується maharṣi, sādhu та yogin — трьох позначень духовних авторитетів, наведених у родовому множини. Текстологічна примітка: SARIT 9.5 передає maharṣīṇāṃ, тоді як GRETIL Avg\_9.5 має maharṣīṣāṃ; це відмінність цифрових читань, а не доведене різночитання рукописів. Тут відтворено IAST чинної редакторської картки.}
\vcom{Рядок б}
\vdeva{दृष्ट्वा निर्वेदम् आपन्नः को न शाम्यति मानवः}
\viast{dṛṣṭvā nirvedam āpannaḥ ko na śāmyati mānavaḥ}
\vtranslit{дріштва нірведам апаннах ко на шам'яті манавах}
\vperek{хто, дійшовши відречення, не втихне?}
\vcom{dṛṣṭvā («побачивши») позначає попередню дію, nirvedam āpannaḥ — людину, яка дійшла стану розчарування чи неприв'язаності, а ko na śāmyati mānavaḥ формує риторичне питання. Вірш не є доказом непридатності кожного вчителя; це інтерпретація зіткнення різних поглядів, а не наказ відмовитися від знання.}
\vcom{Риторичне питання з відповіддю «ніхто». Важливий перелік: `maharṣi`, `sādhu`, `yogin` — три різні авторитети, і текст показує, що **вони не згодні між собою**. Відречення тут походить не від безладдя світу, а від розбіжності авторитетів — це сильніше й конкретніше.}
\vpar{Українська паралель (наша інтерпретація)}{Євген Плужник, «Рання осінь» (1927).}{Наша паралель: у Плужника розгубленість перед безліччю людських поглядів теж приводить до тиші, а не до нового погляду. Наше порівняння — за будовою думки.}
\vfig{ilustratsii/9-5-motion.pdf}{Багато поглядів сходяться до тиші}{Схема наша, не історична ілюстрація.}

\vsec{9.6}
\vdeva{कृत्वा मूर्तिपरिज्ञानं चैतन्यस्य न किं गुरुः । निर्वेदसमतायुक्त्या यस् तारयति संसृतेः ॥}
\viast{kṛtvā mūrtiparijñānaṃ caitanyasya na kiṃ guruḥ | nirvedasamatāyuktyā yas tārayati saṃsṛteḥ ||}
\vtranslit{крітва муртіпаріджнянам чайтаняся на кім ґурух | нірведасаматаюктя яс тараяті самсрітех ||}
\vperek{Хто, пізнавши образ свідомості, — чи не вчитель? — той через відречення, рівність і з'єднаність переправляє з сансари.}
\vcom{Построкове зіставлення: редакторська чернетка, не рецензовано.}
\vcom{Рядок а}
\vdeva{कृत्वा मूर्तिपरिज्ञानं चैतन्यस्य न किं गुरुः}
\viast{kṛtvā mūrtiparijñānaṃ caitanyasya na kiṃ guruḥ}
\vtranslit{крітва муртіпаріджнянам чайтаняся на кім ґурух}
\vperek{Хто, пізнавши образ свідомості, — чи не вчитель? —}
\vcom{kṛtvā mūrtiparijñānaṃ та caitanyasya пов'язують пізнання з образом або виявом свідомості; na kiṃ guruḥ має питальну силу «чи не вчитель?». Синтаксис нетривіальний: перша частина не містить наказу безумовно підкорятися наставникові.}
\vcom{Рядок б}
\vdeva{निर्वेदसमतायुक्त्या यस् तारयति संसृतेः}
\viast{nirvedasamatāyuktyā yas tārayati saṃsṛteḥ}
\vtranslit{нірведасаматаюктя яс тараяті самсрітех}
\vperek{той через відречення, рівність і з'єднаність переправляє з сансари.}
\vcom{nirveda-samatā-yuktyā — форма орудного відмінка: засіб, пов'язаний із неприв'язаністю, рівновагою та способом їх поєднання; yas tārayati saṃsṛteḥ — «хто переправляє із сансари». Образ переправи не вказує конкретної фізичної річки чи історичної події; додавати їх до наукового перекладу не слід.}
\vcom{Єдине місце в книзі, де `guru` названо **позитивно** і з означенням: учитель — той, хто пізнав `caitanya`. Але й тут текст не вимагає послуху: `tārayati` — «переправляє», дієслово руху, а не влади. Три складники (`nirveda`, `samatā`, `yukti`) подані як спосіб дії, а не як чесноти вчителя.}
\vpar{Українська паралель (наша інтерпретація)}{Григорій Сковорода, «Кільце. Дружня розмова про душевний мир».}{Наша паралель: у Сковороди рівність теж є тим, через що людина виходить із кола, у якому її тримали. Наше порівняння — типологічне.}
\vfig{ilustratsii/9-6-witness.pdf}{Переправа з сансари: центр — пізнаний образ свідомості}{Схема наша, не історична ілюстрація.}

\vsec{9.7}
\vdeva{पश्य भूतविकारांस् त्वं भूतमात्रान् यथार्थतः । तत्क्षणाद् बन्धनिर्मुक्तः स्वरूपस्थो भविष्यसि ॥}
\viast{paśya bhūtavikārāṃs tvaṃ bhūtamātrān yathārthataḥ | tatkṣaṇād bandhanirmuktaḥ svarūpastho bhaviṣyasi ||}
\vtranslit{пашя бгутавікарамс твам бгутаматран ятхартхатах | таткшанад бандганірмуктах сварупастхо бгавішясі ||}
\vperek{Дивись на видозміни сущого як на самі лише елементи, як є насправді — і тієї ж миті, вільний від пут, станеш тим, хто стоїть у власній формі.}
\vcom{Построкове зіставлення: редакторська чернетка, не рецензовано.}
\vcom{Рядок а}
\vdeva{पश्य भूतविकारांस् त्वं भूतमात्रान् यथार्थतः}
\viast{paśya bhūtavikārāṃs tvaṃ bhūtamātrān yathārthataḥ}
\vtranslit{пашя бгутавікарамс твам бгутаматран ятхартхатах}
\vperek{Дивись на видозміни сущого як на самі лише елементи, як є насправді —}
\vcom{paśya — наказовий спосіб «дивись»; bhūta-vikārān і bhūta-mātrān співвідносять видозміни й елементи. yathārthataḥ уточнює спосіб сприйняття — «як є насправді». Це заклик до розрізнення в межах філософського тексту, не твердження про сучасну фізику чи медицину.}
\vcom{Рядок б}
\vdeva{तत्क्षणाद् बन्धनिर्मुक्तः स्वरूपस्थो भविष्यसि}
\viast{tatkṣaṇād bandhanirmuktaḥ svarūpastho bhaviṣyasi}
\vtranslit{таткшанад бандганірмуктах сварупастхо бгавішясі}
\vperek{і тієї ж миті, вільний від пут, станеш тим, хто стоїть у власній формі.}
\vcom{tatkṣaṇāt («тієї ж миті») визначає часовий перехід; bandha-nirmuktaḥ («звільнений від пут») і svarūpa-sthaḥ («той, хто перебуває у власній формі») описують очікуваний стан. bhaviṣyasi є майбутнім часом «станеш», а не минулим описом уже досягнутого.}
\vcom{`tatkṣaṇāt` («тієї ж миті») — найсильніше слово вірша: не поступово, не через практику, а негайно. `bhūtamātra` («самі лише елементи») знімає оцінку з речей: вони не погані й не добрі, вони — складники. `svarūpastha` («той, хто стоїть у власній формі») — стан, а не дія; і `bhaviṣyasi` — майбутній час, тому це не опис того, що вже є, а вказівка на те, що станеться.}
\vpar{Українська паралель (наша інтерпретація)}{Володимир Свідзінський, «Медобір» (1975, посмертно).}{Наша паралель: у Свідзінського річ теж розкладається на те, з чого вона є, і саме так стає видимою. Наше порівняння — за образом.}
\vfig{ilustratsii/9-7-thread.pdf}{Видозміни розібрані на елементи, як намистини на нитці}{Схема наша, не історична ілюстрація.}

\vsec{9.8}
\vdeva{वासना एव संसार इति सर्वा विमुञ्च ताः । तत्त्यागो वासनात्यागात् स्थितिरद्य यथा तथा ॥}
\viast{vāsanā eva saṃsāra iti sarvā vimuñca tāḥ | tattyāgo vāsanātyāgāt sthitiradya yathā tathā ||}
\vtranslit{васана ева самсара іті сарва вімуньча тах | таттяґо васанатяґат стхітірадя ятха татха ||}
\vperek{Сансара — це саме усталені схильності, так; відпусти їх усі. Їх відкидання — через відкидання схильностей; стояння в собі — нині, так чи інакше.}
\vcom{Построкове зіставлення: редакторська чернетка, не рецензовано.}
\vcom{Рядок а}
\vdeva{वासना एव संसार इति सर्वा विमुञ्च ताः}
\viast{vāsanā eva saṃsāra iti sarvā vimuñca tāḥ}
\vtranslit{васана ева самсара іті сарва вімуньча тах}
\vperek{Сансара — це саме усталені схильності, так; відпусти їх усі.}
\vcom{vāsanā пов'язано із saṃsāra, а sarvā ... tāḥ указує на множину схильностей, які закликають відпустити; vimuñca — наказова форма. Це філософський контекст звички й прив'язаності, не порада відмовлятися від медичної допомоги, соціальних зв'язків чи відповідальності.}
\vcom{Рядок б}
\vdeva{तत्त्यागो वासनात्यागात् स्थितिरद्य यथा तथा}
\viast{tattyāgo vāsanātyāgāt sthitiradya yathā tathā}
\vtranslit{таттяґо васанатяґат стхітірадя ятха татха}
\vperek{Їх відкидання — через відкидання схильностей; стояння в собі — нині, так чи інакше.}
\vcom{tat-tyāgaḥ та vāsanā-tyāgāt повторюють корінь tyāga («відкидання») у різних граматичних відношеннях; sthitir adya «перебування нині», yathā tathā «так чи інакше». Переклад не повинен вигадувати конкретного ритуалу або гарантувати миттєвий зовнішній результат.}
\vcom{Вірш закінчує главу визначенням, а не обіцянкою: сансара — не світ, а `vāsanā`, усталені схильності. `yathā tathā` («так чи інакше») знімає вимогу до способу: результат не залежить від того, як саме це станеться. Друга половина повторює те саме слово двічі (`tyāga`), і повтор тут не недогляд — він показує, що відкидання схильностей і є відкиданням, а не засобом до нього.}
\vpar{Українська паралель (наша інтерпретація)}{Леся Українка, «Бояриня» (1910).}{Наша паралель: у Лесі Українки схильності, засвоєні звичкою, теж тримають людину міцніше за будь-який наказ. Наше порівняння — типологічне.}
\vfig{ilustratsii/9-8-knot.pdf}{Вузол схильностей і стояння в собі}{Схема наша, не історична ілюстрація.}

